\documentclass[10pt,a4paper]{amsart}
\usepackage[margin=19mm]{geometry}
\usepackage{amsmath,amssymb,mathtools}
\usepackage[scr=rsfs,cal=euler]{mathalfa}
\usepackage{xfrac}
\usepackage[unicode,hidelinks,bookmarks=false]{hyperref}
\newcommand{\ie}{i.e.\ }
\newcommand{\cf}{cf.\ }
\newcommand{\resp}{respectively\ }
\newcommand{\Subsection}{\subsection}
\providecommand{\nobreakdash}{\nobreak\hbox{-}}
\setlength{\emergencystretch}{2em}
\multlinegap0pt
\usepackage{float}
\usepackage{xcolor}
\usepackage{multirow}
\usepackage{mathtools}
\usepackage{amssymb}
\usepackage{booktabs}
\usepackage{tcolorbox}
\usepackage{tikz}
\usepackage{bbm}
\newtheorem{proposition}{Proposition}
\newcommand{\DD}{\mathbb D}
\newcommand{\HH}{\mathbb H}
\newcommand{\RR}{\mathbb R}
\renewcommand{\Re}{\op{Re}}
\newcommand{\cB}{\mathcal B}
\newcommand{\cM}{\mathcal M}
\newcommand{\dave}[1]{#1}\newcommand{\davebegin}{}\newcommand{\daveend}{}\newcommand{\george}[1]{#1}
\newcommand{\eps}{\varepsilon}
\newcommand{\op}{\operatorname}
\newcommand{\ovl}[1]{{\overline{#1}}}
\title{A Spectral Theory of Scalar Volterra Equations}
\author{David Darrow, George Stepaniants}
\date{Source: arXiv:2503.06957v3 (CC BY 4.0)}
\begin{document}
\maketitle

\noindent\emph{Source and licence.} A Spectral Theory of Scalar Volterra Equations. Version arXiv:2503.06957v3 (CC BY 4.0), distributed by arXiv under the Creative Commons Attribution 4.0 International licence, http://creativecommons.org/licenses/by/4.0/, as recorded in the arXiv licence metadata for this exact version and shown on the abstract page https://arxiv.org/abs/2503.06957v3. Full licence notice: \texttt{licenses/arxiv-2503.06957-CC-BY-4.0.txt}.\par\smallskip

\noindent\textit{Editorial scope.} Counted unit: the article's proposition prop:weakcont (source0.tex lines 1810--1812). The article's complete original proof of it is delivered with it (source0.tex lines 1814--1848, one \texttt{proof} environment of 35 lines). No mathematics is written, repaired or re-derived: the statement and the proof are the article's own lines, in the article's own notation, and no retained line is substituted or re-flowed.  No further source-local context is needed: every cross-reference of every retained line resolves inside the two delivered spans.  Nothing was omitted from any retained span, so the package carries no omission record.  No retained line contains a citation, so the wrapper carries no bibliography and no bibliography entry.  No retained line contains a figure environment, an image-inclusion command or drawing code, so no image is delivered and the package declares no asset dependency.  Editorial formatting only, in addition: an AMS \texttt{amsart} preamble that restates the article's own declarations and macros used by the retained lines, namely the environments \texttt{proposition} on separate counters, together with the standard packages those lines need.  The retained environments print in this document as Proposition 1 in order of appearance, whereas the article prints the same items with the numbering of its own full document.  The complete native source of the article as distributed in the arXiv e-print for this exact version is bundled unchanged as \texttt{sources/174474/source0.tex}.
\begin{proposition}[Weak continuity of admissible maps]\label{prop:weakcont}
    If $S$ is admissible, then $\cB_{S}:\cM_+(S^1)\times\RR\to\cM_+(S^1)\times\RR$ is continuous with respect to the product of the weak topology on $\cM_+(S^1)$ and the standard topology on $\RR$.
\end{proposition}

\begin{proof}
    Fix a nonzero $\lambda\in\cM_+(S^1)$ and $\sigma\in\RR$, and suppose $\lambda_n\in\cM_+(S^1)$ converges weakly to $\lambda$ and $\sigma_n$ converges to $\sigma$. For any $r<1$, define the following complex functions on $S^1$:
    \[f_{r,n}(e^{i\theta}) = Q_{\sigma_n}[\lambda_n](re^{i\theta}) = Q[\lambda_n](re^{i\theta}) + i\sigma_n.\]
    Notably, $f_{r,n}$ is smooth, and 
    \[\lim_{n\to\infty} f_{r,n}(e^{i\theta}) = f_r(e^{i\theta})\doteq Q_\sigma[\lambda](re^{i\theta})\]
    pointwise in $S^1$; this follows from the weak convergence of $\lambda_n$, as the Cauchy kernel is smooth and uniformly bounded along each fixed $r$. Suppose $\|\lambda\|_{S^1} = M>0$, and fix $N\geq 1$ such that $\|\lambda_{n}\|_{S^1}\leq 2M$ for all $n\geq N$. We then know that $\partial_\theta f_{r,n}$ is uniformly bounded in $n$, since
    \[|\partial_\theta f_{r,n}(e^{i\theta})| = \frac{1}{2\pi}\left|\int_0^{2\pi}\frac{1+ire^{i(\theta-\theta')}}{1-ire^{i(\theta-\theta')}}\,d\lambda_n(\theta')\right| \leq \frac{2M}{2\pi}\,\frac{1+r}{1-r},\]
    \dave{and so a standard argument shows that $\lim_{n\to\infty} f_{r,n}=f_r$
    uniformly, for fixed $r$}. Fix a neighborhood $U\supset f_r(S^1)$ in $\HH_+$; by increasing $N$, we can guarantee that 
    \[f_{r,n}(S^1)\subset U\]
    for $n\geq N$. However, $S$ is smooth on $\ovl{U}$, so in particular, it is uniformly Lipschitz on $U$; as such,
    \[\lim_{n\to\infty} S\circ f_{r,n} = S\circ f_r\]
    uniformly, for fixed $r$.

    Let $(\mu,\xi) = \cB_{S}[\lambda,\sigma]$ and $(\mu_n,\xi_n) = \cB_{S,\sigma}[\lambda_n,\sigma_n]$. By applying the convergence of $S\circ f_{r,n}$ to the case $r=0$, we see that $\xi_n\to\xi$ and $\|\lambda_n\|\to\|\lambda\|$ as $n\to\infty$. As such, if $\|\mu\|_{S^1} = M'$, we can increase $N$ to ensure that $\|\mu_n\|_{S^1}\leq 2M'$ for $n\geq N$. Next, define the following measures in $\cM_+(S^1)$:
    \[\mu_{r,n} = (2\pi)^{-1} (S^{\Re}\circ f_{r,n})(e^{i\theta})\,d\theta,\qquad \mu_r = (2\pi)^{-1} (S^{\Re}\circ f_r)(e^{i\theta})\,d\theta,\]
    and fix a bounded, continuous function $g:S^1\to\RR$. Since $P[g]$ is continuous on the compact set $\ovl{\DD}$, it is necessarily uniformly continuous. For any $\eps>0$, then, we can fix $r_\eps<1$ such that
    \[\sup\nolimits_\theta\left|g(e^{i\theta}) - P[g](re^{i\theta})\right|<\eps\]
    for $r\geq r_\eps$. Define $\widetilde{g}(e^{i\theta})=g(e^{-i\theta})$. Then we find
    \begin{align*}
        \int g\,d\mu_{r,n} &= (2\pi)^{-1}\int_0^{2\pi} g(\theta) P[\mu_n](re^{i\theta})\,d\theta\\
        &= (2\pi)^{-1}\int_0^{2\pi} g(\theta) \int_0^{2\pi} \Re\left(\frac{1+re^{i(\theta-\theta')}}{1-re^{i(\theta-\theta')}}\right)\,d\mu_n(\theta')\,d\theta\\
        &= \int_0^{2\pi} P[\widetilde{g}](re^{-i\theta'})\,d\mu_n(\theta')\\
        &= \int_0^{2\pi} P[g](re^{i\theta'})\,d\mu_n(\theta'),
    \end{align*}
    and similarly for $\mu_r$ and $\mu$; this implies
    \begin{align*}\left|\int g\,d(\mu_{n} - \mu) \right| &\leq \left|\int g\,d(\mu_{n,r_\eps} - \mu_n) \right| + \left|\int g\,d(\mu_{n,r_\eps}-\mu_{r_\eps}) \right| + \left|\int g \,d(\mu_{r_\eps} - \mu) \right|\\
    &\leq \left|\int_0^{2\pi} \left(g(e^{i\theta})-P[g](r_\eps e^{i\theta})\right)\,d\mu_n(\theta) \right| + \left|\int g\,d(\mu_{n,r_\eps}-\mu_{r_\eps}) \right| \\
    &\qquad+ \left|\int_0^{2\pi} \left(g(e^{i\theta})-P[g](r_\eps e^{i\theta})\right)\,d\mu(\theta) \right|\\
    &\leq \left|\int g\,d(\mu_{n,r_\eps}-\mu_{r_\eps}) \right| + 2\eps M'.
    \end{align*}
    Since $\eps$ was arbitrary and $\mu_{n,r_\eps}\rightharpoonup\mu_{r_\eps}$ weakly, we deduce that
    \[\int g\,d\mu_n\to \int g\,d\mu\]
    for any bounded, continuous function $g$. The proposition follows.
\end{proof}

\end{document}
